\documentclass{subfiles}
\begin{document}
    Our algorithm assume the existence of a function \lstinline{Random(a, b)},
        that for any integers \(a, b\) gives a uniformely distributed random integer \(r\)
        in the range \([a, b]\). We also assume \lstinline{Random(a, b)} to run in \(\ofO{1}\).

    Let \(\sigma = \iprod{\sigma_{1}, \ldots, \sigma_{m}}\) be a stream of distint tokens.
        We say that the \emph{rank} of \(\sigma_{i}\) is the number of elements in \(\sigma\)
        smaller the \(\sigma_{i}\), plus 1:
        \[
            \rank(\sigma_{i}) = 1 + \abs{\set{ \sigma_{i}' \in \sigma}[\sigma_{i}' < \sigma_{i}]}.
        \]
        We call median an item whose rank is either \(\ceil{\sfrac{m + 1}{2}}\) or 
        \(\floor{\sfrac{m + 1}{2}}\). The algorithm we describe tries to solve the following problem:
        given a stream \(\sigma\) of \(m\) distinct tokens over the universe \(U\),
        compute the median of \(\sigma\).

    Let us consider the following naive solution: let us set \(r = \) \lstinline{Random(1, m)},
        then our algorithm returns \(\sigma_{r}\) as the median element.
        The algorithm takes \(\ofO{\log(n + m)}\), 
        although the length of the stream needs to be known a priori.

    This simple algorithm is far from a good algorithm. 
        Let us observe in fact that, if \(X = \rank(\sigma_{i})\) is a random variable,
        one has that 
        \[
            \E[X]{} = \sum_{ i = 1 }^{m} i \cdot \log\frac{1}{m} = \frac{m + 1}{2}.
        \]
        Although the expected value of the reported item is what we would hope for,
        this is not enough. Indeed, if we would always report the largest (or the smallest) 
        item in the stream, the expectation would still be the same but the item 
        would be far from the median.

    Thus, we need a way to determine how close the reported item is to the median.
        We do so, by considering the following quantity:
        \[
            \floor*{\left\lparen\frac{1}{2} - \varepsilon\right\rparen (m + 1)}
                \le \rank(\sigma_{i})
                \le \ceil*{\left\lparen\frac{1}{2} + \varepsilon\right\rparen (m + 1)}
        \]
        where \(0 \le \varepsilon \le \sfrac{1}{2}\). An element that satisfies the
        above equation is called an \(\varepsilon\)-approximate median.

    To finish of, we now introduce a technique that significantly boost our success rate.
        The idea is extreamly simple: we run the algorithm \(k\) times,
        and report the median of the \(k\) execution as the median of the whole stream.
        This simple trick, allows to define an algorithm that takes \(\ofO{k \log(n + m)}\)
        bits of storage and reports an approximate median with a probability of 
        \(1 - 2(\sfrac{e}{4})^{\sfrac{k}{4}}\).
\end{document}
